\documentclass[a4paper,11pt]{article}
\usepackage[utf8]{inputenc}
\usepackage[french]{babel}
\usepackage[margin=2.5cm]{geometry}
\usepackage{amsmath}
\usepackage{amssymb}
\usepackage{graphicx}
\usepackage{xcolor}
\usepackage{enumitem}

\title{\textbf{Examen UEC3 -- Architectures et ingénierie\\des systèmes de transmission}}
\author{Décembre 2020}
\date{}

\begin{document}

\maketitle

\section*{Consignes}
\begin{itemize}
    \item Documents de cours/PC/BE et calculatrice autorisés
    \item Les parties (A, B) sont indépendantes et sont à rédiger sur des copies séparées
    \item Vous déposerez vos copies sur Moodle, dans l'espace de dépôt prévu, sous forme de scan de votre document manuscrit ou via un document électronique (.pdf)
\end{itemize}

\section{Architectures et ingénierie des systèmes de transmission RF (45 mn)}

On se propose d'étudier la faisabilité d'une liaison radio aux fréquences 'TéraHertz' (f > 300 GHz) dédiée aux transmissions courtes portées à très hauts débits. Les applications visées vont des systèmes optiques avec des ponts radio, à la connexion large bande de serveurs de données en passant par la télévision HDTV. La figure 1 illustre un exemple d'application.

\begin{figure}[h]
    \centering
    \textit{Fig. 1 : lien optique avec extension par un lien radio TéraHertz (THz)}
\end{figure}

Du fait de la très forte atténuation à ces fréquences liée à l'atmosphère ainsi qu'à la pluie, seules certaines bandes (« fenêtres ») sont utilisables (voir annexe A1). Au-delà de 300 GHz, les fenêtres utilisables, définies à partir de la figure A1.a en annexe, et pour lesquelles l'atténuation linéique est inférieure à 100 dB/km, sont données dans le tableau I ci-dessous. Elles offrent potentiellement de très grandes largeurs de bandes.

Compte-tenu des composants disponibles à ces fréquences, seule la fenêtre I est retenue dans le cadre de cette étude et pour un lien fixe.

\subsection{Bilan de liaison}

\begin{enumerate}[label=\alph*.]
    \item Dans les conditions de propagation en espace libre, donner l'expression générale de l'atténuation en espace libre $A_{EL}$ en fonction de la distance $d$ et de la fréquence $f$.
    
    \item Calculer $A_{EL}$ pour $d = 1$ m, 10 m et 100 m.
    
    \item Quelle est l'atténuation linéique maximale due à l'atmosphère et à la pluie, avec un taux de précipitation $R = 35$ mm/h, pour la bande I (cf annexe A1) ?
    
    \item Donner dans un tableau l'atténuation totale $A_T$ (en dB) subie par le signal pour les 3 distances considérées.
\end{enumerate}

\subsection{Architecture}

À 340 GHz, l'architecture du récepteur THz est relativement simple et est composée des éléments représentés à la figure 2 ci-dessous. Le signal à la fréquence intermédiaire de 50 GHz est traité dans un récepteur spécifique. La bande passante du signal transmis, c'est-à-dire la largeur d'un canal, est $B = 10$ GHz.

\begin{figure}[h]
    \centering
    \textit{Fig. 2 : Architecture du récepteur THz}
\end{figure}

Les caractéristiques des composants sont données dans le tableau suivant :

\textit{[Tableau des caractéristiques fourni dans le sujet]}

\begin{enumerate}[label=\alph*., start=5]
    \item Préciser le type d'architecture utilisé et donner la fréquence de l'oscillateur local.
    
    \item Calculer le facteur de bruit total $F_R$ (en dB) du récepteur
    
    \item Une antenne de type cornet circulaire est connectée à l'accès E. Donner l'expression de la puissance de bruit totale $N_R$ en fonction de la température équivalente de bruit totale du récepteur $T_R$ et calculer $N_R$ sachant que $T_R = 2900$ K.
    
    \item On note $P_{TX}$ (dBm) la puissance à l'émission. On suppose que les antennes d'émission et de réception sont identiques et de gain $G_A$. Donner l'expression de la puissance reçue $P_{RX}$ (dBm) en fonction de $A_T$, $G_A$ et $P_{TX}$.
    
    \item Donner l'expression du rapport signal-à-bruit (SNR) en dB en fonction de $A_T$, $G_A$, $P_{TX}$ et $N_R$.
    
    \item On prend $P_{TX} = 1$ mW. Connaissant $N_R$ (question g), donner l'expression du gain des antennes $G_A$ en fonction de $A_T$ et du SNR requis.
    
    \item Calculer le gain $G_A$ nécessaire pour garantir un SNR de 15 dB pour les 3 distances (1, 10 et 100 m). Commentaires sur le diamètre des antennes résultant.
\end{enumerate}

\subsection*{Annexe A1 : Atténuations linéiques (atmosphère et pluie)}

\textit{[Figures A1.a et A1.b fournies montrant l'atténuation en fonction de la fréquence]}

\section{Ingénierie d'une liaison amplifiée (45 mn)}

En utilisant les caractéristiques des équipements fournies en annexe, vous devez proposer une ingénierie de liaison WDM amplifiée en bande C répondant aux exigences suivantes :

\begin{enumerate}
    \item La capacité totale utile transmise Cap doit être maximisée.
\end{enumerate}

\subsection*{Travail demandé}

\textit{[Tableau à compléter avec les paramètres du système]}

\subsection*{Annexe : Caractéristiques des équipements}

\subsubsection*{Ligne optique}
\textit{[Caractéristiques de la fibre et de la ligne]}

\subsubsection*{Transpondeur (caractéristiques de la partie réception)}
\textit{[Caractéristiques des transpondeurs]}

\subsubsection*{Amplificateur optique}
\textit{[Caractéristiques des EDFA]}

\end{document}
